\section{预备知识}
\label{sec:preliminaries-zh}

本节定义论文中使用的核心概念。
异构图记为 $\mathcal{G}=(\mathcal{V},\mathcal{E},\mathbf{X},\phi,\psi)$，其中 $\mathcal{V}$ 和 $\mathcal{E}$ 分别表示节点集和边集，$\mathbf{X}$ 为节点特征矩阵，$\phi$ 和 $\psi$ 分别给出节点类型和边类型。
当图中包含多种节点类型或边类型时，该图称为异构图。

元路径是在异构图模式上定义的复合关系。
它描述了节点类型和关系类型按照某种顺序组合形成的语义路径，例如学术网络中的 author--paper--author 表示共同写作关系，author--paper--venue--paper--author 表示通过发表 venue 连接作者。
给定一条元路径，可以得到某个节点的元路径邻居集合，不同元路径会产生不同语义邻域。

基于元路径的异构图学习通常先根据多条元路径构造多个语义视图，再在每个视图中聚合邻居信息，最后通过语义级融合得到节点表示。
本文中的拓扑视图正是建立在这种语义聚合范式之上。

结构对抗攻击是指在扰动预算约束下修改边集合，得到受攻击图 $\widetilde{\mathcal{G}}$。
攻击者希望在尽量少改动结构的情况下破坏节点分类性能。
本文关注结构扰动下的鲁棒异构图学习，即在受攻击拓扑上仍学习可靠的目标类型节点表示。
